\section{Related Work}
\label{sec:related}

Three bodies of work approach the question of source identity effects in SFT, but none closes
the gap we address.

\subsection{Domain Mixture Optimization for Language Models}

The influence of training data composition on model performance is well established for
pretraining. DoReMi \cite{Xie2023DoReMi} optimizes domain weights via Group Distributionally
Robust Optimization, demonstrating that the proportion of text from different domains
significantly affects downstream performance. Chameleon \cite{Xie2025Chameleon} uses leverage-score
weighting for heterogeneous corpora in both pretraining and finetuning settings. DomainPilot
\cite{Zhang2026DomainPilot} applies domain-level loss monitoring to optimize SFT data mixture, achieving
+3.8\% on LiveCodeBench.

However, these approaches optimize mixture proportions from a fixed pool of diverse sources.
They do not isolate the causal effect of training on a single source dataset identity. The
question ``what happens if you train only on HumanEval problems versus only on MBPP problems?''
remains unasked.

\subsection{Data Composition for Code SFT}

Code SFT papers universally report training data compositions but do not ablate source
identity. WizardCoder \cite{Luo2023WizardCoder} and OctoPack \cite{Muennighoff2023OctoPack} demonstrate
that augmenting data quality via Evol-Instruct or instruction formatting improves HumanEval
pass@1 --- but they vary data quality and style simultaneously with source, making causal
attribution impossible. DeepSeek-Coder \cite{Guo2024DeepSeekCoder} reports fixed compositions without
source isolation. Lv et al.\ \cite{Lv2025DataEfficient} show saturation in code SFT data selection but use a
curated multi-source dataset. Chen et al.\ \cite{Chen2024UnlockSFT} ablate atomic versus synthetic source types ---
the closest existing work, but not source dataset identity with a cross-benchmark transfer
matrix.

\subsection{Distributional Alignment as a Predictor}

Zhang et al.\ \cite{Zhang2025GRAPE} (GRAPE) show that selecting instruction-tuning data by distribution
alignment outperforms 3$\times$ more data --- the most direct precedent for our mechanistic claim.
However, GRAPE operates on general instruction-following, not code-specific SFT, without
code-specialized embeddings or permutation testing.

Our contribution extends this line by providing the first permutation-tested evidence
(CodeBERT Spearman $\rho$=1.0, p=0.042, n=10,000) that embedding-space alignment between SFT
source and test benchmark governs pass@1 rank order for code generation.
